\usepackage{hyperref}
\usepackage{anyfontsize} % kill warnings about font sizes
\usepackage{amsmath}     % Mathematics FTW!
\usepackage[capitalise]{cleveref}
%\usepackage{amsthm}      % Warning, I've only proven it correct :).
%\usepackage{amssymb}     % needed for \mathbb and mathcal
%\usepackage{bbm}
\usepackage[english]{babel}
%\usepackage{amsbsy}      % for \boldsymbol
%\usepackage{stmaryrd}   % provides some PLy symbols, like Scott Brackets
\usepackage{mathpartir}  % for \mathpar
\usepackage{color}
\usepackage{braket}      % for \set and \braket
\usepackage{mathtools}   % for \bigtimes
\usepackage{thmtools}    % to duplicate theorems in the appendix
%\usepackage{marvosym} % for deniarius (gradual label)
%\usepackage{wasysym}     % for \smiley
%\usepackage{balance}     % for the references
\usepackage{xspace}
%\usepackage{mathrsfs} % for mathscr font
\usepackage{proof}
% \usepackage{yhmath}
\usepackage{rotating}
\usepackage[Symbol]{upgreek} % upright greek letters
\usepackage{tikz-cd}
\usepackage{xcolor}
\usepackage{changepage}
\usepackage{centernot}
\usepackage{comment}
\usepackage{listings}
\usepackage{mdframed}
\usepackage{centernot}
\usepackage{mathtools}
\usepackage{multirow}
\usepackage{xparse}
\definecolor{darkblue}{rgb}{0.0, 0.0, 0.55}
\definecolor{darkgreen}{rgb}{0.0, 0.55, 0.13}
\definecolor{darkred}{rgb}{0.55, 0.0, 0.0}
\definecolor{darkviolet}{rgb}{0.58, 0.0, 0.83}
\definecolor{lightblue}{rgb}{0.68, 0.85, 0.9}
\usepackage{lstcoq}
% Literate Agda
% \usepackage{agda}
\usepackage[all]{xy}
\usepackage{stmaryrd}
\usepackage{subcaption,verbatim}
\usepackage[utf8x]{inputenc}
% \usepackage{agda-syntax}
\usepackage[outline]{contour}


\DeclareMathAlphabet{\mathbbm}{U}{bbm}{m}{n}

%% hack to get upright greek letters
\renewcommand{\Omega}{\Upomega}
\renewcommand{\Pi}{\Uppi}
\renewcommand{\Sigma}{\Upsigma}
\renewcommand{\Gamma}{\Upgamma}
\renewcommand{\Delta}{\Updelta}

\makeatletter
% \@ifpackageloaded{stix}{%
% }{%
%   \DeclareFontEncoding{LS2}{}{\noaccents@}
%   \DeclareFontSubstitution{LS2}{stix}{m}{n}
%   \DeclareSymbolFont{stix@largesymbols}{LS2}{stixex}{m}{n}
%   \SetSymbolFont{stix@largesymbols}{bold}{LS2}{stixex}{b}{n}
%   \DeclareMathDelimiter{\lBrace}{\mathopen} {stix@largesymbols}{"E8}%
%                                             {stix@largesymbols}{"0E}
%   \DeclareMathDelimiter{\rBrace}{\mathclose}{stix@largesymbols}{"E9}%
%                                             {stix@largesymbols}{"0F}
% }
\makeatother
\DeclareMathAlphabet{\mathcal}{OMS}{cmsy}{m}{n} % Reset mathcal font because ACM version sucks

% Comment form:  requires amssymb and color

\newcommand{\mynote}[2]
    {{\color{red} \fbox{\bfseries\sffamily\scriptsize#1}
    {\small$\blacktriangleright$\textsf{\emph{#2}}$\blacktriangleleft$}}~}

\definecolor{propcolor}{HTML}{3F7D31}
\definecolor{mypurple}{HTML}{5B069D}
\newcommand{\myproposal}[2]
    {{\color{propcolor} \fbox{\bfseries\sffamily\scriptsize#1}
    {\small$\blacktriangleright$\textsf{\emph{#2}}$\blacktriangleleft$}}~}

\newcommand{\nt}[1]{\mynote{NT}{#1}}
\newcommand{\lp}[1]{\mynote{LP}{#1}}
\newcommand{\jc}[1]{\mynote{JC}{#1}}
\newcommand{\tw}[1]{\mynote{TW}{#1}}
\newcommand{\pmp}[1]{\mynote{PMP}{#1}}

% Grey Box
\definecolor{lightgray}{gray}{0.90}
\newcommand{\Gbox}[1]{\colorbox{lightgray}{$#1$}}

% text
\newcommand{\ie}{\emph{i.e.,}\xspace}
\newcommand{\etal}{\emph{
    et al.}\xspace}
\newcommand{\eg}{\emph{e.g.,}\xspace}

\newcommand{\nat}{\mathbb{N}}
\newcommand{\natP}{\mathbb{N}_\mathbb{P}}
\newcommand{\bool}{\mathsf{Bool}}

\newcommand{\slogan}[1]{\begin{center}``\emph{#1}''\end{center}}

\newcommand{\agdahtml}{https://htmlpreview.github.io/?https://github.com/CoqHott/logrel-mltt/blob/impredicativity-cast-compute-refl/html/}
\newcommand{\refAgdaRoot}[1]{{\small{\href{\agdahtml/#1.html}{\emph{[#1.agda]}}}}}
\newcommand{\refAgdaDef}[1]{{\small{\href{\agdahtml/Definition.#1.html}{\emph{[#1.agda]}}}}}
\newcommand{\refAgda}[2]{{\small{\href{\agdahtml/Definition.#1.#2.html}{\emph{[#2.agda]}}}}}
\newcommand{\refCoq}[1]{{\color{darkblue}\emph{[#1]}}}
% \newcommand{\vrefCoq}[2]{\href{https://github.com/TheoWinterhalter/template-coq/tree/rewrite-rules/pcuic/theories/#1}{\emph{[#2]}}}
\newcommand{\vrefCoq}[2]{\href{https://github.com/TheoWinterhalter/template-coq/tree/popl21-artifact/pcuic/theories/#1}{\emph{[#2]}}}


% \newcommand{\Coq}{${\mathrm{Coq}}$\xspace}
\newcommand{\MLTT}{\ensuremath{\mathsf{MLTT}}\xspace}
\newcommand{\sMLTT}{\ensuremath{\mathsf{sMLTT}}\xspace}
\newcommand{\ETT}{\ensuremath{\mathsf{ETT}}\xspace}
\newcommand{\HTS}{\ensuremath{\mathsf{HTS}}\xspace}
\newcommand{\SetoidTT}{\ensuremath{\mathsf{TT}^\mathrm{obs}}\xspace}
\newcommand{\SetoidCC}{\ensuremath{\mathsf{CC}^\mathrm{obs}}\xspace}
\newcommand{\SetoidCIC}{\ensuremath{\mathsf{CIC}^\mathrm{obs}}\xspace}
\newcommand{\CIC}{\ensuremath{\mathsf{CIC}}\xspace}

\newcommand{\UIP}{${\mathrm{UIP}}$\xspace}

\newcommand{\bbfamily}{\fontencoding{U}\fontfamily{bbold}\selectfont}
% \DeclareMathAlphabet{\mathbbold}{U}{bbold}{m}{n}
\newcommand{\Two}{\omega}
%% \newcommand{\Two}{{\color{white}\contour{blga}}
\newcommand{\independent}{\emptyset}

\newcommand{\prop}{\prim{Prop}}
\newcommand{\Univ}{\oblset{Univ}}
\newcommand{\hProp}{\oblset{hProp}}
\def\Coq{\textsc{Coq}\xspace}
\def\Agda{\textsc{Agda}\xspace}
\def\Lean{\textsc{Lean}\xspace}
\def\Idris{\textsc{Idris}\xspace}
% \newcommand{\Agda}{${\mathrm{Agda}}$\xspace}
% \newcommand{\Lean}{${\mathrm{Lean}}$\xspace}
% \newcommand{\Cedille}{${\mathrm{Cedille}}$\xspace}
\newcommand{\CICU}{${\mathrm{CIC}}_u$\xspace}
\newcommand{\RTT}{${\mathrm{RTT}}$\xspace}
\def\MetaCoq{\textsc{MetaCoq}\xspace}

\newcommand{\subst}[3][x]{#2[{#1}:={#3}]}

\newcommand\dashPar{\vdash}%_{p}}

\newcommand{\sizeName}{\mathtt{size}}
\newcommand{\size}[1]{\sizeName(#1)}

\newcommand{\wfctx}[2][]{#1 \vdash #2}
\newcommand{\eqctx}[2]{\vdash #1 \equiv #2}
\newcommand{\tytm}[4][]{#1 #2 \vdash #3 : #4}
\newcommand{\tytmPara}[4][]{#1 #2 \dashPar #3 : #4}
\newcommand{\tm}[2]{#1 : #2}
\newcommand{\tytms}{\tytm}
\newcommand{\eqtm}[5][]{#1 #2 \vdash #3 \equiv #4 : #5}
\newcommand{\nered}[4]{#1 \vdash #2 \Rightarrow_{ne} #3 : #4}
\newcommand{\neconv}[4]{#1 \vdash #2 \cong_{ne} #3 : #4}
\newcommand{\neconvRed}[4]{#1 \vdash #2 \cong_{ne}^\downarrow #3 : #4}
\newcommand{\whnfconv}[4]{#1 \vdash #2 \cong #3 : #4}
\newcommand{\whnfconvRed}[4]{#1 \vdash #2 \cong^\downarrow #3 : #4}
\newcommand{\red}[5][]{#1 #2 \vdash #3 \Rightarrow #4 : #5}
\newcommand{\redT}[5][]{#1 #2 \vdash #3 \Rightarrow^* #4 : #5}
\newcommand{\eqtmmultiline}[6][]{#1 #2 \vdash \begin{array}{l} #3
                                                \equiv \\
                                                #4 \end{array} : #5 : #6}
\newcommand{\redmultiline}[5][]{#1 #2 \vdash \begin{array}{l} #3 \Rightarrow \\ #4 \end{array} : #5}
\newcommand{\redmultilineExt}[6][]{#1 #2 \vdash \begin{array}{l} #3 \Rightarrow \\ #4 \\ #5 \end{array} : #6}

\newcommand{\eqproj}[3]{#1_\epsilon^{#2}~#3}
\newcommand{\propext}{\sProp{}_\mathrm{ext}}
\newcommand{\funext}{\Pi_\mathrm{ext}}

\newcommand{\tytyR}[4]{#2 \Vdash_{#1} #3 : #4}
\newcommand{\tytmR}[5]{#2 \Vdash_{#1} #3 : #4 : #5}
\newcommand{\ctxV}[1]{\Vdash^v #1}
\newcommand{\tytyV}[4]{#2 \Vdash^v_{#1} #3 : #4}
\newcommand{\tytyS}[3]{#1 \Vdash^s #2 : #3}
\newcommand{\tytmV}[5]{#2 \Vdash^v_{#1} #3 : #4 : #5}
\newcommand{\nattmR}[2]{#1 \Vdash_{\mathbb{N}\mathsf{t}}\ {#2}}
\newcommand{\nateqR}[3]{#1 \Vdash_{\mathbb{N}\mathsf{t}=}\ {#2} \equiv {#3}}
\newcommand{\quotmR}[2]{#1 \Vdash_Q\ {#2}}
\newcommand{\quoeqR}[3]{#1 \Vdash_Q\ {#2} \equiv {#3}}
\newcommand{\idtmR}[2]{#1 \Vdash_\mathrm{Id}\ {#2}}
\newcommand{\ideqR}[3]{#1 \Vdash_\mathrm{Id}\ {#2} \equiv {#3}}
\newcommand{\eqtyR}[5]{#2 \Vdash_{#1} #3 \equiv #4 : #5}
\newcommand{\eqtyS}[4]{#1 \Vdash^s #2 \equiv #3 : #4}
\newcommand{\eqtmR}[6]{#2 \Vdash_{#1} #3 \equiv #4 : #5 : #6}
\newcommand{\eqtmV}[6]{#2 \Vdash^v_{#1} #3 \equiv #4 : #5 : #6}
\newcommand{\wk}[3]{{#1}}% : {#2} \le {#3}}
\newcommand{\wkrho}{(\rho : \Delta \sqsubseteq \Gamma)}
\newcommand{\wkrhoshort}{\wkrho}

\newcommand{\metatm}[2]{{#1} :: {#2}}
\newcommand{\metaforall}[2]{\forall {#1},\ {#2}}
\newcommand{\metaimply}[2]{{#1}\ \to\ {#2}}

\newcommand{\emptyctx}{\bullet}
\newcommand{\extctx}[3][x]{#2 , \tm{#1}{#3}}
\newcommand{\inctx}[4][x]{\tmannotated{#1}{#3}{#4}{} \in {#2}}

\newcommand{\Type}{\mathsf{Type}}
\newcommand{\SProp}{\mathsf{SProp}}
\newcommand{\Prop}{\mathsf{Prop}}

\newcommand{\forded}[1]{#1 _\mathbb{F}}
\newcommand{\eqpi}[2][]{#2 _\epsilon^{#1}}
\newcommand{\appi}[1]{#1 _{ap}}
\newcommand{\Depsum}[3][x]{\Sigma (\tm{#1}{#2}).\, {#3}}
\newcommand{\Depfun}[3][x]{\Pi (\tm{#1}{#2}).\, {#3}}
\newcommand{\Fun}[2]{{#1} \to {#2}}
\newcommand{\Sig}[3][x]{\Sigma (\tm{#1}{#2}).\, {#3}}
\newcommand{\Prod}[2]{{#1} \land {#2}}
\newcommand{\Exists}[3][x]{(\tm{#1}{#2}) \ \& \ {#3}}
\newcommand{\varExists}[3][x]{\exists (\tm{#1}{#2})\ .\ {#3}}
\newcommand{\ExistsUnique}[3][x]{\exists! (\tm{#1}{#2})\ .\ {#3}}
\newcommand{\Path}[3][]{#2 \sim_{#1} #3}
\newcommand{\eqInd}[3]{\mathrm{eq}_{#1} \, #2 \, #3}
\newcommand{\Empty}{\bot}
\newcommand{\botP}{\bot_{\mathbb{P}}}
\newcommand{\Unit}{\top}
\newcommand{\Quo}[2]{{#1}/{#2}}
\newcommand{\Id}[3]{\metaop{Id}(#1,#2,#3)}
\newcommand{\Boxt}[1]{\square{#1}}
\newcommand{\Squash}[1]{\| #1 \|}
\newcommand{\CubicalPath}{\metaop{Path}}
\newcommand{\List}[1]{\mathtt{list} \, #1}

\newcommand\metaop[1]{{\color{RoyalBlue}{\mathrm{#1}}}}

\newcommand{\piRel}[2]{\mathcal{R} ( #1 , #2)}
\newcommand{\univRel}[2]{\mathcal{A} ( #1 , #2)}

\newcommand{\lstop}[1]{\textrm{\coqeN{#1}}\xspace}
\newcommand{\lstops}[1]{\textrm{\lstinline[language=Coq, basicstyle=\small,escapechar=*]{#1}}\xspace}

\newcommand{\emptyrec}[2]{\mathtt{\color{darkblue} \bot\mathrm{-} elim}~{#1}~{#2}}
\newcommand{\unit}{*}
\newcommand{\zero}{{\color{darkblue} 0}}
\newcommand{\sucName}{\mathtt{\color{darkblue} S}}
\newcommand{\suc}[1]{\sucName\ #1}
\newcommand{\natrec}[4]{\mathtt{\color{darkblue} \mathbb{N}\mathrm{-} elim}~{#1}~{#2}~{#3}~{#4}}
\newcommand{\lam}[3][x]{\lambda (\tm{#1}{#2}).\, #3}
\newcommand{\pair}[2]{\langle #1 , #2 \rangle }
\newcommand{\fst}[1]{\metaop{fst}({#1})}
\newcommand{\snd}[1]{\metaop{snd}({#1})}
% \newcommand{\refl}[2]{\lstop{refl}~#1}
\newcommand{\refl}[2]{\mathtt{\color{darkblue} refl}~#1}
\newcommand{\reflI}[2]{\mathtt{\color{darkblue} refl}~#1}
\newcommand{\sym}[1]{{#1}^{-1}}
\newcommand{\transitivity}[2]{{#1}\cdot{#2}}
\newcommand{\ap}[2]{\lstop{ap}~{#1}~{#2}}
\newcommand{\castName}{\lstop{cast}}
% \newcommand{\cast}[4]{{\castName~#1~#2~#3~#4}}
\newcommand{\cast}[4]{{\mathtt{\color{darkblue} cast}~#1~#2~#3~#4}}
\newcommand{\castI}[4]{{\mathtt{\color{darkblue} cast}~#1~#2~#3~#4}}
\newcommand{\transportName}{\lstop{transp}}
% \newcommand{\transport}[6]{\transportName~#1~#2~#3~#4~#5~#6}
\newcommand{\transport}[6]{\mathtt{\color{darkblue} transport}~#1~#2~#3~#4~#5~#6}
\newcommand{\transportI}[6]{\mathtt{\color{darkblue} transport}~#1~#2~#3~#4~#5~#6}
\newcommand{\castreflName}{\metaop{castrefl}}
\newcommand{\castrefl}[2]{\castreflName(#1,#2)}
\newcommand{\quo}[1]{\pi(#1)}
\newcommand{\quorec}[4]{\metaop{Q\mathrm{-}elim}(#1,#2,#3,#4)}
\newcommand{\id}[1]{\metaop{Idrefl}(#1)}
\newcommand{\idpath}[1]{\metaop{Idpath}(#1)}
\newcommand{\J}[6]{\metaop{J}(#1,#2,#3,#4,#5,#6)}
\newcommand{\eqJ}[5]{\metaop{eq_J}(#1,#2,#3,#4,#5)}
\newcommand{\boxt}[1]{\diamond{#1}}
\newcommand{\unboxt}[1]{\metaop{\Box\mathrm{-elim}}(#1)}
\newcommand{\squash}[1]{| #1 |}
\newcommand{\unsquash}[3]{\metaop{\mathrm{S}\mathrm{-elim}}(#1,#2,#3)}
\newcommand{\nil}{\mathtt{nil}}
\newcommand{\cons}[2]{\mathtt{cons} \, #1 \, #2}

\newcommand\head[1]{\mathop{\text{hd}}#1}
\newcommand\whnf[1]{\mathop{\text{whnf}}#1}
\newcommand\neutral[1]{\mathop{\text{neutral}} #1}
\newcommand\postype[1]{\mathop{\text{posType}}#1}



\newcommand{\tmannotated}[3]{#1 : #2 : #3}%{#1 :^{#3} #2}
\newcommand{\extctxannotated}[4][x]{#2 , \tmannotated{#1}{#3}{#4}{}}
\newcommand{\tytmannotated}[5][]{#1 #2 \vdash \tmannotated{#3}{#4}{#5}}
\newcommand{\tytmParaannotated}[5][]{#1 #2 \dashPar \tmannotated{#3}{#4}{#5}}
\newcommand{\eqtmannotated}[6][]{#1 #2 \vdash #3 \equiv #4 : #5 : #6}
\newcommand{\algoeqtmannotated}[6][]{#1 #2 \vdash #3 \cong^\downarrow #4 : #5 : #6}
\newcommand{\eqtmParaannotated}[6][]{#1 #2 \dashPar #3 \equiv #4 : #5 : #6}
\newcommand{\redannotated}[5][]{#1 #2 \vdash #3 \Rightarrow #4 : #5}
\newcommand{\Depfunannotated}[6][x]{\Pi_{}^{#5, #6} (\tm{#1}{#2}).\, {#3}}
\newcommand{\Depsumannotated}[6][x]{\Sigma_{#4,#5}^{#6} (\tm{#1}{#2}).\, {#3}}
\newcommand{\Existsannotated}[5][x]{(\tm{#1}{#2}) \ \& \ {#3}}

\newcommand{\Con}{\mathsf{Con}}
\newcommand{\cwfty}[1][]{\mathsf{Ty}_{#1}}
\newcommand{\cwftm}[1][]{\mathsf{tm}_{#1}}
\newcommand{\cwfprop}[1][]{\mathsf{Pr}_{#1}}
\newcommand{\cwfpr}[1][]{\mathsf{pr}_{#1}}

\newcommand{\Setoid}{\mathsf{Setoid}}
\newcommand{\metasProp}{\mathsf{Prop}}
\newcommand{\metanat}{\boldsymbol{\omega}}
\newcommand{\metazero}{0}
\newcommand{\metasuc}[1]{\mathsf{S}\ #1}
\newcommand{\metaempty}{\bot}
\newcommand{\metaunit}{\top}
\newcommand{\metarefl}{\mathsf{refl}}

\renewcommand{\ne}{N}

\newcommand{\Context}{\AgdaData{Context}}
\newcommand{\Term}{\AgdaData{Term}}
\newcommand{\Sort}{\AgdaData{Sort}}

\newcommand{\VU}{{\textcolor{AgdaInductiveConstructor} \Vdash}_{\AgdaCon{U}}}
\newcommand{\VOmega}{{\textcolor{AgdaInductiveConstructor} \Vdash}_{\textcolor{AgdaInductiveConstructor} \Omega}}
\newcommand{\Vempty}{{\textcolor{AgdaInductiveConstructor} \Vdash}_{\textcolor{AgdaInductiveConstructor} \bot}}
\newcommand{\Vpi}{{\textcolor{AgdaInductiveConstructor} \Vdash}_{\textcolor{AgdaInductiveConstructor} \Pi}}
\newcommand{\Vforall}{{\textcolor{AgdaInductiveConstructor} \Vdash}_{\textcolor{AgdaInductiveConstructor} \forall}}
\newcommand{\Vexists}{{\textcolor{AgdaInductiveConstructor} \Vdash}_{\textcolor{AgdaInductiveConstructor} \&}}
\newcommand{\Vnat}{{\textcolor{AgdaInductiveConstructor} \Vdash}_{\AgdaCon{N}}}
\newcommand{\Vne}{{\textcolor{AgdaInductiveConstructor} \Vdash}_{\AgdaCon{ne}}}
\newcommand{\Vemb}{{\textcolor{AgdaInductiveConstructor} \Vdash}_{\AgdaCon{emb}}}
\newcommand{\VX}{{\textcolor{AgdaInductiveConstructor} \Vdash}_{\AgdaCon{X}}}

\newcommand{\RU}{\AgdaCon{R}_{\AgdaCon{U}}}
\newcommand{\ROmega}{\AgdaCon{R}_{\textcolor{AgdaInductiveConstructor} \Omega}}
\newcommand{\Rempty}{\AgdaCon{R}_{\textcolor{AgdaInductiveConstructor} \bot}}
\newcommand{\Rpi}{\AgdaCon{R}_{\textcolor{AgdaInductiveConstructor} \Pi}}
\newcommand{\Rforall}{\AgdaCon{R}_{\textcolor{AgdaInductiveConstructor} \forall}}
\newcommand{\Rexists}{\AgdaCon{R}_{\textcolor{AgdaInductiveConstructor} \exists}}
\newcommand{\Rnat}{\AgdaCon{R}_{\AgdaCon{N}}}
\newcommand{\Rne}{\AgdaCon{R}_{\AgdaCon{ne}}}
\newcommand{\Remb}{\AgdaCon{R}_{\AgdaCon{emb}}}

\newcommand{\sV}{\mathbf{V}}
\newcommand{\sU}{\mathbf{U}}
\newcommand{\Upred}[1][i]{\mathsf{code}_{#1}}
\newcommand{\sType}{\mathbf{s}}
\newcommand{\ssProp}{\Omega}
\newcommand{\sFun}[2]{#1 \to_s #2}
\newcommand{\sEq}[2]{#1 = #2}
\newcommand{\el}{\mathsf{el}}
\newcommand{\eleq}{\mathsf{el\text{-}eq}}
\newcommand{\val}{\mathsf{val}}
\newcommand{\valeq}{\mathsf{val\text{-}eq}}
\newcommand{\codeemb}{\mathsf{c}_\mathsf{emb}}
\newcommand{\codenat}{\mathsf{c}_\mathsf{N}}
\newcommand{\codepi}{\mathsf{c}_{\Pi}}
\newcommand{\codepiuu}{\mathsf{c}_{\Pi\mathsf{UU}}}
\newcommand{\codepisu}{\mathsf{c}_{\Pi\Omega\mathsf{U}}}
\newcommand{\codepius}{\mathsf{c}_{\Pi\mathsf{U}\Omega}}
\newcommand{\codepiss}{\mathsf{c}_{\Pi\Omega\Omega}}
\newcommand{\codeex}{\mathsf{c}_{\exists}}
\newcommand{\codeU}{\mathsf{c}_\mathsf{U}}
\newcommand{\codesProp}{\mathsf{c}_{\sProp}}
\newcommand{\codeInd}{\mathsf{c}_\mathsf{ind}}
\newcommand{\codeempty}{\mathsf{c}_\bot}
\newcommand{\codequo}{\mathsf{c}_Q}

\newcommand{\mctx}[1]{\llbracket\ #1\ \rrbracket}
\newcommand{\mty}[2][\Gamma]{\llbracket {#2} \rrbracket_{#1}}
\newcommand{\mtm}[3][\Gamma]{[\tm{#2}{#3}]_{#1}}

\newcommand{\bnfis}{::=}
\newcommand{\bnfin}{\in}
\newcommand{\sep}{|}
\newcommand{\qsep}{\ | \ }

% Automatically insert parentheses
\newcommand{\maybeparens}[2]{%
  \ifthenelse{\isempty{#2}}{#1}{%  if argument is empty
  \ifthenelse{\equal{#2}{(}}{#1#2}{%  or (, then do not parenthesize
  #1(#2)}}} % otherwise, parenthesize
\newcommand{\tdom}{\mathrm{dom}}
\newcommand{\dom}{\maybeparens{\tdom}}
\newcommand{\tPV}{\mathrm{PV}}
\newcommand{\PV}{\maybeparens{\tPV}}
\newcommand{\tFV}{\mathrm{FV}}
\newcommand{\FV}{\maybeparens{\tFV}}

\newcommand{\joinsub}{\uplus}%{\|}
\newcommand{\Joinsub}{\biguplus}%{\|}

\newcommand{\mathsc}[1]{\ensuremath{\textsc{#1}}}
\newcommand{\rulename}[1]{{\footnotesize\mathsc{#1}}}
\newcommand{\nru}[3]{\dfrac{\begin{array}{@{}c@{}}#1\end{array}}{\begin{array}{@{}c@{}}#2\end{array}}\ifthenelse{\isempty{#3}}{}{\ \rulename{#3}}}
\newcommand{\ru}[2]{\nru{#1}{#2}{}}

\newcommand{\myref}[2]{\hyperref[#2]{#1~\ref*{#2}}}
% Referencing figures
\newcommand{\figref}[1]{\myref{Fig.}{fig:#1}}
\newcommand{\secref}[1]{\myref{Sect.}{sec:#1}}
\newcommand{\defref}[1]{\myref{Definition}{def:#1}}
\newcommand{\lemref}[1]{\myref{Lemma}{lem:#1}}
\newcommand{\thmref}[1]{\myref{Theorem}{thm:#1}}
\newcommand{\eqnref}[1]{\hyperref[#1]{(\ref{#1})}}
\newcommand{\exref}[1]{\myref{Example}{ex:#1}}



% Agda highlighting macros
%% \newcommand{\data}[1]{{\AgdaDatatype{#1}}}
%% \newcommand{\record}[1]{{\AgdaRecord{#1}}}
%% \newcommand{\func}[1]{{\AgdaFunction{#1}}}
%% \newcommand{\prim}[1]{{\AgdaPrimitive{#1}}}
%% \newcommand{\primty}[1]{{\AgdaPrimitiveType{#1}}}
%% \newcommand{\sset}[1]{{\prim{Set}_{#1}}}
%% \newcommand{\var}[1]{{\AgdaBound{#1}}}
%% \newcommand{\keyw}[1]{{\AgdaKeyword{#1}}}
%% \newcommand{\con}[1]{{\AgdaInductiveConstructor{#1}}}
%% \newcommand{\field}[1]{{\AgdaField{#1}}}
%% \newcommand{\proj}[1]{\field{.#1}}
%% \newcommand{\pragma}[2]{\texttt{\{-\# } \AgdaKeyword{\texttt{#1}}\ #2 \texttt{ \#-\}}}

%% \newcommand{\Id}[1]{\mathrel{{\AgdaPrimitiveType{\equiv}}_{#1}}}
%% \newcommand{\Leq}{\mathrel{\AgdaDatatype{\leq}}}
%% \newcommand{\Simeq}{\mathrel{\AgdaRecord{\simeq}}}
%% \newcommand{\sumtype}{\mathrel{\AgdaDatatype{\uplus}}}
%% \newcommand{\prodtype}{\mathrel{\AgdaDatatype{\times}}}

%% \newcommand{\pat}[1]{\AgdaKeyword{?}#1}

%% \newcommand{\Nat}{{\color{AgdaDatatype}\mathbb{N}}}
%% \newcommand{\myNat}{\func{N}}
%% \newcommand{\vect}[2]{\func{Vec} \ #1 \ #2}
%% \newcommand{\myZero}{\func{zero}}
%% \newcommand{\Sone}{\func{$S^1$}}
%% \newcommand{\base}{\func{base}}
%% \newcommand{\Loop}{\func{loop}}
%% \newcommand{\mySucc}{\func{suc}}
%% \newcommand{\myNatRec}{\func{elim}}%_{\myNat}}
%% \newcommand{\myCatch}{\func{catch}}
%% \newcommand{\myRaise}{{\func{raise}}}
%% \newcommand{\Exc}{\AgdaPostulate{$\mathbb{E}$}}
%% \newcommand{\Int}{\data{\mathbb{Z}}}
%% \newcommand{\toptype}{{\color{AgdaDatatype}\top}}
%% \newcommand{\bottomtype}{\data{\bot}}
%% \newcommand{\sigmatype}{\data{\Sigma}}
%% \newcommand{\Bool}{{\color{AgdaDatatype}\mathbb{B}}}
%% \newcommand{\myMax}{\AgdaFun{max}}

%% \newcommand{\Fin}{\data{Fin}}
%% \newcommand{\fzero}{\con{fzero}}
%% \newcommand{\fsuc}{\con{fsuc}}
%% \newcommand{\lezero}{\con{lz}}
%% \newcommand{\lesuc}{\con{ls}}
%% \newcommand{\refl}{\con{refl}}
%% \newcommand{\reflbar}{\overline{\con{refl}}}
%% \newcommand{\zero}{\con{zero}}
%% \newcommand{\suc}{\con{suc}}
%% \newcommand{\true}{\con{true}}
%% \newcommand{\false}{\con{false}}
%% \newcommand{\J}{\prim{J}}
%% \newcommand{\K}{\prim{K}}

%% \newcommand{\zero}{{\color{AgdaDatatype}\mathsf{0}}}
%% \newcommand{\succ}{{\color{AgdaDatatype}\mathsf{S}}}
%% \newcommand{\ty}[1][]{{\color{AgdaDatatype}\square}_{#1}}
%% \newcommand{\natind}{\Nat_{\mathsf{ind}}}
%% \newcommand{\natind}{\myNatRec}

%% \newcommand{\rwRule}[4]{#1 \ctxrw \cmd{#2}{#3} \rwred #4}

%% \newcommand{\cmd}[2]{\langle #1 \mid #2 \rangle}
%% \newcommand{\cmdEval}[2]{\mathrm{eval}\, ( #1 \mid #2 )}
%% \newcommand{\zip}[2]{\mathsf{zip}(#1,\ #2)}
%% \newcommand{\zip}[2]{\mathsf{zip}(#1 \mid #2)}
%% \newcommand{\ctxrw}{\rightthreetimes}
%% \newcommand{\ctxrw}{\ltimes}
%% \newcommand{\confluent}{\downdownarrows}
%% \newcommand{\force}[1]{{\Downarrow}(#1)}
%% \newcommand{\patt}[1]{\AgdaKeyword{#1}}
%% \newcommand{\patt}[1]{{#1}}
%% \newcommand{\pattinst}[3]{#2 [ #3 ]_{#1}}
%% \newcommand{\force}[1]{{\patt{\Downarrow}}(#1)}
%% \newcommand{\erp}[1]{#1^{-}}
%% \newcommand\cored{\leftsquigarrow}
%% \newcommand{\rectype}[1]{\mathfrak{#1}}
%% \newcommand{\recmk}[2]{\lBrace{#2}\rBrace}
%% \newcommand{\recproj}[3]{{#2}._{#1} #3}
%% \newcommand{\stkproj}[2]{{\diamond}._{#1} #2}
%% \newcommand{\stkproj}[2]{{\diamond}.#2}
%% \newcommand{\stkapply}[1]{\diamond\ #1}
%% \newcommand{\inst}{\Rightarrow}
%% \newcommand{\patstk}{\mathbb{S}}
%% \newcommand{\patframe}{\mathbb{F}}

%% \newcommand{\reds}{\red^\star}
%% \newcommand{\coreds}{\cored^\star}

%% \newcommand\pos[2]{{% we make the whole thing an ordinary symbol
%%   \left.\kern-\nulldelimiterspace % automatically resize the bar with \right
%%   #1 % the function
%%   % \vphantom{\big|} % pretend it's a little taller at normal size
%%   \right|_{#2} % this is the delimiter
%%   }}


%% \newcommand\optRedName{\rho}
%% \newcommand\optRed[1]{\optRedName(#1)}
%% \def\paraRed{\Rrightarrow}
%% %\def\red{\rightarrow}
%% \newcommand\closure[1]{#1^*}
%% \def\Cred{\closure{\red}}

%% \newcommand\mgu[2]{\textit{MGU}(#1,#2)}

%% \newcommand{\RewContext}{\mathcal{R}}
%% \newcommand{\RewOrder}{\succ}

%% \newcommand\whnfName[1]{\mathtt{whnf}_{#1}}
%% \newcommand\whnf[2]{\whnfName{#1}(#2)}


%%%%%%%%%%%%%%%%%%%%%%%%%%%%%%%%%%%%%%%%%%%%%%%%%%%%%%%%%%%%%%%%%%%%%%%%%%%%%%%%%%%%%%%%%%%%%%%%%%%%%%%%%%%%%%%%%%%%%
%Reference to inference rules, inspired by https://tex.stackexchange.com/questions/340788/cross-referencing-inference-rules

\makeatletter
\let\originferrule\inferrule
\DeclareDocumentCommand \inferrule { s O {} m m }{%
	\IfBooleanTF{#1}%
	{%
		\mpr@inferstar[#2]{#3}{#4}%
	}{%
		\mpr@inferrule[#2]{#3}{#4}%
	}%
  % Modification of the SO code: we reuse the main label
	\IfValueT{#2}%
	{%
		\my@name@inferrule{#2}%
	}%
}
\NewDocumentCommand \my@name@inferrule { m }{%
	\def\@currentlabelname{\textsc{#1}}%
}

% \DeclareDocumentCommand \redrule {o m m}{%
% 	\IfValueTF{#1}
% 	{\textsc{#1}: \; #2 \redCCIC #3 \my@name@inferrule{#1}}
% 	{#2 \redCCIC #3}
% }

\newcommand*{\ilabel}[1]{%
  % \edef\@currentlabel{#1}% Set target label
  \phantomsection% Correct hyper reference link
  \label{#1}% Print and store label
}
\makeatother
